\chapter{40}

\indentedblock{Damals. Fribourg, fünf.\\
Louis ist zappelig.

«Morgen gehen wir in eine heilige Messe,» hat sein Vater gesagt.

Endlich ist es soweit. Es ist Sonntagmorgen.

Sie sitzen zu dritt in der schweren, schmalen Kirchenbank. Der Vater in
der Mitte, Louis links, Lucia rechts neben dem Vater.

Louis weiss, dass sie still sitzen und ruhig sein sollen. Mit offenem
Mund versucht er die Worte des Pfarrers zu verstehen. Er hat noch
niemanden auf diese Weise reden hören. Der Mann singt mehr, als dass er
spricht. Das macht Louis mächtig Eindruck.
Ebenso die Menschen um ihn herum, die alle ganz ernst in die gleiche
Richtung schauen. Und zur gleichen Zeit genau das Gleiche tun. Sein
Vater ist immer ein bisschen zu spät dran. Egal ob sie nun auf einmal
aufstehen sollen, sich hinknien oder setzen müssen, Louis sieht das. Er
versucht möglichst schnell, es ihm gleichzutun. Lucia hat es leichter.
Der Vater zieht sie gleich selbst in jede neue Position.

Dann plötzlich erschrickt Louis. Mit einem Donner setzt gewaltige
Orgelmusik ein. Dazu erklingen tiefe, männliche Stimmen. Er kann nicht einordnen, was der
Gesang, der von hinten und weit oben kommt, mit ihm macht. Er schwebt
auf ihn nieder, klingt düster, schwer und drückt.

Louis schaut verloren um sich, bis seine Augen bei diesem 
jungen Mann stehen bleiben, den er von Vaters Erzählungen kennt. Das muss Jesus sein. 
Er starrt ihn an. Er ist an das Holz genagelt. Eine rote Flüssigkeit fliesst über seine Hände und Füsse. Weint er? 
Warum hängt er eigentlich da oben, wenn er doch so lieb ist. Später wird er den Vater 
fragen, warum der Mann nur eine zerrissene Unterhose trägt.

Louis verbietet sich, zu blinzeln. Vielleicht wird Jesus plötzlich
hinunter kommen. Von diesem Kreuz, direkt zu ihm. Wenn er kurz
wegschaut. Immer wieder hält er den Atem an, bis er kaum mehr Luft
bekommt.

Louis kann sich nicht erklären, was dem armen Mann passiert ist. Er tut
ihm leid. Vielleicht muss er auch einmal so sterben. Maman geht nie mit
ihnen in solche Gottesdienste. Louis weiss nicht, warum. Sie ist
einfach dagegen. Er sehnt sich nach ihr. Sie würde ihm jetzt die Hand
halten. Ganz bestimmt.

Panisch rutscht er ganz nach vorne an die Kante der Bank, bis seine
Füsse den durchgetretenen Holzboden berühren. Er drückt seine Schuhe
darauf und presst die Hände ineinander. Den Vater hat er an diese fremde
Stimmung längst verloren. Auch seine Schwester ist unerreichbar. Sie
ist, angelehnt an den Vater, einfach eingeschlafen.}

\sectionbreak

Louis ging in die Küche hoch und goss sich Wasser in die Trinkflasche. Was hatte ihm Manuela mit diesem Musiktipp nur eingebrockt. Die längst vergessene Kirchenszene war so plötzlich wie detailliert in ihm abgelaufen, als befände er sich gerade wieder auf dieser Bank.
Eigentlich hatte sich sein Vater mit der \emph{heiligen Episode} selbst überboten. Louis musste knapp im Kindergartenalter gewesen sein, als es angefangen hatte. Der Vater hatte von bedingungsloser Liebe geschwärmt. Er hatte Louis dabei an sich gedrückt gehalten und ihm immer wieder diese Geschichte erzählt. Atemlos fasziniert hatte Louis ihm zugehört.
Und Maman sich regelmässig entfernt. Der junge Mann war ihm gleich sympathisch gewesen. Jesus. Der Name hatte sich je länger desto tiefer in ihm eingebrannt. Wenn der Vater erzählte, fühlte Louis innige Verbundenheit. Er gehörte ihm dann ganz alleine. Dem ersten waren ein paar wenige weitere Gottesdienstbesuche gefolgt. Nur, so beeindruckend wie Jesus aufgetaucht war, so plötzlich war er auch wieder verschwunden. Und mit ihm auch ihre Kirchenbesuche.

Louis setzte sich an den kleinen Tisch und stütze den Kopf auf. Weitere stillgelegte Erinnerungen
taten sich als verwirrende Szenen auf. Vaters Aktivitäten mit ihm hatten
sich gejagt. Und der Ablauf sich ein aufs neue Mal wiederholt.

Die Karatekleidung hatte Maman eines Tages im Keller verschwinden
lassen. Louis liebte sie noch heute dafür. Er hatte sich vor den
Konkurrenten arg gefürchtet.
Kaum überstanden, hatte ihm der Vater eines Tages von Motorrädern
geschwärmt. Louis sollte eines bekommen. Der Vater hatte von echten
Motorradrennen für Kinder erzählt. Vor Eifer hatte er sich beim Sprechen
dauernd verhaspelt. Und als er ihm Bilder von kleinen Maschinen zeigte,
die echten Motorrädern in nichts nachstanden, fielen Louis fast die
Augen aus dem Gesicht. Seine anfängliche Faszination hatte sich aber
bald in Abwehr verwandelt. Vaters Kunde mit dem kleinen Sohn hatte den
Schrecken befeuert.

«Dem zeigen wir, wer den besseren Sohn hat. Den schlägst du leicht,
Louis», hatte der Vater zu ihm gesagt, als hätten sie bereits gewonnen.

Damit hatte er ihn in ein Nest übermütiger Racker geworfen und seine
Erwartungen hatten von Rennen zu Rennen Louis’ anfänglichen Stolz
erstickt. Einzig Vaters uneingeschränkte Zuwendung während der
Trainingsfahrten hatten ihn die Nachmittage überleben lassen. Bis er
stürzte und mit gebrochenem Arm am Boden lag. 

Das Bild liess Louis zusammenzucken und eine unbändige Wut kam hoch. Er massierte sich die Schläfen. Wie es Maman nur gelungen war, ihn von diesem Kindersport zu erlösen? Denn das hatte sie. 
 
Louis schüttelte den Kopf, hob die Info-Mappe hoch und stand auf. 
Er wollte diese Erinnerungen nicht; doch die Abwehr blieb chancenlos. Offenbar stand die verriegelte Tür zu seiner Kindheit  unterdessen weit offen. Louis stöhnte.
Was war sein Vater für ein Idiot. Der Gedanke liess seine Hände den Ordner wie ein Schutzschild an die Brust pressen.
Wie wenig Vaters «Projekte» seinen Träumen entsprochen hatten, und woher der nur immer wieder diese
Ideen genommen hatte? 
Einige Jahre später dann die ersehnte Wende. Rückblickend hatten ihn die
Fussballschuhe gerettet.
Vater hatte Louis im Club anmelden wollen. Er sollte ein grosser
Fussballer werden. Warum, war ihm unerklärlich geblieben. Er hatte sich doch alle Mühe
gegeben, ihm seine Argumente gegen diesen Sport darzulegen. Und dann
hatte der doch tatsächlich diese neongelben Nockenschuhe unter seiner
Jacke hervorgezaubert. An einem schulfreien Nachmittag. Sie ihm
hingestreckt, gejubelt und gesagt, am Abend beginne sein erstes
Training. Und Louis hatte geschrien. Aus lauter Verzweiflung.

Das erste Mal hatte er ihm vehement widersprochen.

Das erste Mal hatte er seinen Vater als wutschnaubenden Gegner
empfunden.

Und auch das erste Mal hatte sich seine Mutter ganz direkt zwischen ihn und
den Vater gestellt. Der hatte die Schuhe gepackt, sie mit Wucht in
den Abfalleimer geworfen und die Wohnung mit einem Türknallen verlassen.

\sectionbreak

Wie weit unten Louis die Ereignisse schon versteckt hielt, über all die
Jahre. Ob er Vaters Aktionen je verstehen würde?

Wo sein Vater wohl gerade war? Die Frage hatte er sich schon sehr lange
nicht mehr gestellt.
